%%% Hlavní soubor. Zde se definují základní parametry a odkazuje se na ostatní části. %%%

%% Verze pro jednostranný tisk:
% Okraje: levý 40mm, pravý 25mm, horní a dolní 25mm
% (ale pozor, LaTeX si sám přidává 1in)
\documentclass[hidelinks,12pt,a4paper]{report}
\setlength\textwidth{145mm}
\setlength\textheight{247mm}
\setlength\oddsidemargin{15mm}
\setlength\evensidemargin{15mm}
\setlength\topmargin{0mm}
\setlength\headsep{0mm}
\setlength\headheight{0mm}
% \openright zařídí, aby následující text začínal na pravé straně knihy
\let\openright=\clearpage

%% Pokud tiskneme oboustranně:
% \documentclass[12pt,a4paper,twoside,openright]{report}
% \setlength\textwidth{145mm}
% \setlength\textheight{247mm}
% \setlength\oddsidemargin{14.2mm}
% \setlength\evensidemargin{0mm}
% \setlength\topmargin{0mm}
% \setlength\headsep{0mm}
% \setlength\headheight{0mm}
% \let\openright=\cleardoublepage

%% Vytváříme PDF/A-2u
\usepackage[a-2u]{pdfx}

%% Přepneme na českou sazbu a fonty Latin Modern
\usepackage[czech]{babel}
\usepackage{lmodern}
\usepackage[T1]{fontenc}
\usepackage{textcomp}

%% Použité kódování znaků: obvykle latin2, cp1250 nebo utf8:
\usepackage[utf8]{inputenc}

%%% Další užitečné balíčky (jsou součástí běžných distribucí LaTeXu)
\usepackage{amsmath}        % rozšíření pro sazbu matematiky
\usepackage{amsfonts}       % matematické fonty
\usepackage{amsthm}         % sazba vět, definic apod.
\usepackage{bbding}         % balíček s nejrůznějšími symboly
			    % (čtverečky, hvězdičky, tužtičky, nůžtičky, ...)
\usepackage{bm}             % tučné symboly (příkaz \bm)
\usepackage{graphicx}       % vkládání obrázků
\usepackage{fancyvrb}       % vylepšené prostředí pro strojové písmo
\usepackage{indentfirst}    % zavede odsazení 1. odstavce kapitoly
\usepackage{natbib}         % zajištuje možnost odkazovat na literaturu
			    % stylem AUTOR (ROK), resp. AUTOR [ČÍSLO]
\usepackage[nottoc]{tocbibind} % zajistí přidání seznamu literatury,
                            % obrázků a tabulek do obsahu
\usepackage{icomma}         % inteligetní čárka v matematickém módu
\usepackage{dcolumn}        % lepší zarovnání sloupců v tabulkách
\usepackage{booktabs}       % lepší vodorovné linky v tabulkách
\usepackage{paralist}       % lepší enumerate a itemize
\usepackage{xcolor}         % barevná sazba
\usepackage{multicol}

\usepackage{rotating}
%\usepackage{fontspec}
\usepackage{tipa}
\usepackage{enumitem}
\usepackage{makecell}
\usepackage{alltt}
\usepackage{refcount}
\usepackage{tabularx}
\usepackage{caption}

%\usepackage{etoolbox}
%\makeatletter
%\patchcmd{\chapter}{\if@openright\cleardoublepage\else\clearpage\fi}{}{}{}
%\makeatother


%%% Údaje o práci

% Název práce v jazyce práce (přesně podle zadání)
\def\NazevPrace{Svěcení kněžstva v~Církvi československé husitské a církvi Římskokatolické}

% Název práce v angličtině
\def\NazevPraceEN{Ordination of Clergy in the Czechoslovak Hussite Church and in the Roman Catholic Church}

% Jméno autora
\def\AutorPrace{Mgr. Jan Oldřich Krůza}

% Rok odevzdání
\def\RokOdevzdani{2021}

% Název katedry nebo ústavu, kde byla práce oficiálně zadána
% (dle Organizační struktury MFF UK, případně plný název pracoviště mimo MFF)
\def\Katedra{Katedra historické teologie a církevních dějin}
\def\KatedraEN{Department of Historical Theology and Church History}

% Jedná se o katedru (department) nebo o ústav (institute)?
\def\TypPracoviste{Katedra}
\def\TypPracovisteEN{Department}

% Vedoucí práce: Jméno a příjmení s~tituly
\def\Vedouci{JUDr. Jiří Sedlák}

% Pracoviště vedoucího (opět dle Organizační struktury MFF)
\def\KatedraVedouciho{Katedra historické teologie a církevních dějin}
\def\KatedraVedoucihoEN{Department of Historical Theology and Church History}

% Studijní program a obor
\def\StudijniProgram{husitská teologie}
\def\StudijniObor{HT NK}

% Nepovinné poděkování (vedoucímu práce, konzultantovi, tomu, kdo
% zapůjčil software, literaturu apod.)
%\def\Podekovani{%
%Děkuji
%}

% Abstrakt (doporučený rozsah cca 80-200 slov; nejedná se o zadání práce)
\def\Abstrakt{%
Tato seminární práce porovnává institut svěcení kněžstva v~Církvi československé
husitské a v~církvi Římskokatolické z~pohledu člena první jmenované církve.
Krátce reflektuje nad kněžstvím jako takovým, odlišnostech v~konceptech a
v~požadavcích a stručně zhodnotí příčiny a důsledky těchto rozdílů.
}
\def\AbstraktEN{%
This thesis compares the institution of clergy ordination in the Czechoslovak
Hussite Chuch and in the Roman Catholic Church from the point of view of a
member of the former one. It shortly reflects upon priesthood as such, upon
conceptual differences and in requirements for ordination and it briefly
evaluates causes and effects of said differences.
}

% 3 až 5 klíčových slov (doporučeno), každé uzavřeno ve složených závorkách
\def\KlicovaSlova{%
{svěcení}, {kněžství}, {Církev československá husitská}, {Římskokatolická církev}
}
\def\KlicovaSlovaEN{%
{ordination}, {clergy}, {Czechoslovak Hussite Church}, {Roman Catholic Church}
}

%% Balíček hyperref, kterým jdou vyrábět klikací odkazy v PDF,
%% ale hlavně ho používáme k uložení metadat do PDF (včetně obsahu).
%% Většinu nastavítek přednastaví balíček pdfx.
\hypersetup{unicode}
\hypersetup{breaklinks=true}

%% Definice různých užitečných maker (viz popis uvnitř souboru)
%%% Tento soubor obsahuje definice různých užitečných maker a prostředí %%%
%%% Další makra připisujte sem, ať nepřekáží v ostatních souborech.     %%%

%%% Drobné úpravy stylu

% Tato makra přesvědčují mírně ošklivým trikem LaTeX, aby hlavičky kapitol
% sázel příčetněji a nevynechával nad nimi spoustu místa. Směle ignorujte.
\makeatletter
\def\@makechapterhead#1{
  {\parindent \z@ \raggedright \normalfont
   \Huge\bfseries \thechapter. #1
   \par\nobreak
   \vskip 20\p@
}}
\def\@makeschapterhead#1{
  {\parindent \z@ \raggedright \normalfont
   \Huge\bfseries #1
   \par\nobreak
   \vskip 20\p@
}}
\makeatother

% Toto makro definuje kapitolu, která není očíslovaná, ale je uvedena v obsahu.
\def\chapwithtoc#1{
\chapter*{#1}
\addcontentsline{toc}{chapter}{#1}
}

% Trochu volnější nastavení dělení slov, než je default.
\lefthyphenmin=2
\righthyphenmin=2

% Zapne černé "slimáky" na koncích řádků, které přetekly, abychom si
% jich lépe všimli.
%\overfullrule=1mm

%%% Makra pro definice, věty, tvrzení, příklady, ... (vyžaduje baliček amsthm)

\theoremstyle{plain}
\newtheorem{veta}{Věta}
\newtheorem{lemma}[veta]{Lemma}
\newtheorem{tvrz}[veta]{Tvrzení}

\theoremstyle{plain}
\newtheorem{definice}{Definice}

\theoremstyle{remark}
\newtheorem*{dusl}{Důsledek}
\newtheorem*{pozn}{Poznámka}
\newtheorem*{prikl}{Příklad}

%%% Prostředí pro důkazy

\newenvironment{dukaz}{
  \par\medskip\noindent
  \textit{Důkaz}.
}{
\newline
\rightline{$\qedsymbol$}
}

%%% Prostředí pro sazbu kódu, případně vstupu/výstupu počítačových
%%% programů. (Vyžaduje balíček fancyvrb -- fancy verbatim.)

\DefineVerbatimEnvironment{code}{Verbatim}{fontsize=\small, frame=single}

%%% Prostor reálných, resp. přirozených čísel
\newcommand{\R}{\mathbb{R}}
\newcommand{\N}{\mathbb{N}}

%%% Užitečné operátory pro statistiku a pravděpodobnost
\DeclareMathOperator{\pr}{\textsf{P}}
\DeclareMathOperator{\E}{\textsf{E}\,}
\DeclareMathOperator{\var}{\textrm{var}}
\DeclareMathOperator{\sd}{\textrm{sd}}

%%% Příkaz pro transpozici vektoru/matice
\newcommand{\T}[1]{#1^\top}

%%% Vychytávky pro matematiku
\newcommand{\goto}{\rightarrow}
\newcommand{\gotop}{\stackrel{P}{\longrightarrow}}
\newcommand{\maon}[1]{o(n^{#1})}
\newcommand{\abs}[1]{\left|{#1}\right|}
\newcommand{\dint}{\int_0^\tau\!\!\int_0^\tau}
\newcommand{\isqr}[1]{\frac{1}{\sqrt{#1}}}

%%% Vychytávky pro tabulky
\newcommand{\pulrad}[1]{\raisebox{1.5ex}[0pt]{#1}}
\newcommand{\mc}[1]{\multicolumn{1}{c}{#1}}

\newcommand{\ridv}[3]{\href{http://radio.makon.cz/zaznam/#1\#ts=#3}{\texttt{#1} od pozice #2}}
\newcommand{\rid}[3]{\href{http://radio.makon.cz/zaznam/#1\#ts=#3}{\texttt{#1} #2}}
\newcommand{\rref}[3]{\footnote{\rid{#1}{#2}{#3}}}


%% Titulní strana a různé povinné informační strany
\begin{document}
%%% Titulní strana práce a další povinné informační strany

%%% Titulní strana práce

\pagestyle{empty}
\hypersetup{pageanchor=false}

\begin{center}

{\bf\Large UNIVERZITA KARLOVA}

\vfill

{\Large HUSITSKÁ TEOLOGICKÁ FAKULTA}

\vfill

\centerline{\mbox{\includegraphics[width=66mm]{rc/logo-htf.jpg}}}

\vfill

{\bfseries\NazevPrace}

\vfill

{\bfseries\NazevPraceEN}

\vfill

Diplomová práce

\vspace{15mm}

\begin{multicols}{2}
Vedoucí práce:\\
\Vedouci\\
\columnbreak
Autor:\\
\AutorPrace\\
\end{multicols}

\vfill

% Zde doplňte rok
Praha \RokOdevzdani

\end{center}

\newpage

%%% Strana s čestným prohlášením k disertační práci

\openright
\hypersetup{pageanchor=true}
\pagestyle{plain}
\pagenumbering{roman}
\vglue 0pt plus 1fill

\noindent
Prohlašuji, že jsem tuto seminární práci vypracoval samostatně a výhradně
s~použitím citovaných pramenů, literatury a dalších odborných zdrojů.

\medskip\noindent
Beru na~vědomí, že se na moji práci vztahují práva a povinnosti vyplývající
ze zákona č. 121/2000 Sb., autorského zákona v~platném znění, zejména skutečnost,
že Univerzita Karlova má právo na~uzavření licenční smlouvy o~užití této
práce jako školního díla podle §60 odst. 1 autorského zákona.

\vspace{10mm}

\hbox{\hbox to 0.5\hsize{%
%V \hbox to 6em{\dotfill} dne \hbox to 6em{\dotfill}
V Praze dne \today
\hss}\hbox to 0.5\hsize{\dotfill\quad}}
\smallskip
\hbox{\hbox to 0.5\hsize{}\hbox to 0.5\hsize{\hfil Podpis autora\hfil}}

\vspace{20mm}
\newpage

%%% Poděkování

\openright

\noindent
%\Podekovani

\newpage

%%% Povinná informační strana disertační práce

\openright

\vbox to 0.5\vsize{
\setlength\parindent{0mm}
\setlength\parskip{5mm}

{\bf Abstrakt:\\}
\Abstrakt

{\bf Klíčová slova:\\}
\KlicovaSlova

\vfill

{\bf Abstract:\\}
\AbstraktEN

{\bf Keywords:\\}
\KlicovaSlovaEN

\vss}

\newpage



\openright


%%% Strana s automaticky generovaným obsahem disertační práce

\tableofcontents

\newpage
\pagenumbering{arabic}
\setcounter{page}{1}
%%% Jednotlivé kapitoly práce jsou pro přehlednost uloženy v samostatných souborech

\chapter{Úvod}
\label{kap:uvod}

\section{Vymezení pojmu svátost}
\label{div:vymezeni}

Co je vlastně svátost? Jak tomu v~humanitních oborech bývá zvykem, pojem není
ostře a bezesporně ohraničen. Teologie je navíc obor, který se zabývá věcmi
člověka přesahujícími, ba tento svět přesahujícími. O co více tedy je obtížné
definovat pojmy, jako to umí např. matematika? Přesná definice je možná,
pohybujeme-li se v~rámci dobře definovaného modelu. Hůře je možná, pohybujeme-li
se v~rámci popisu reálného světa. A ještě hůře je možná, jestliže hranice světa
překračujeme.

S~rostoucí obtížností podat definici ale roste i relevance a obsažnost pojmu.
Např. v~hypotetickém modelu ve formě počítačové hry mohu svátost zcela přesně
definovat jako nějaký úkon nebo předmět, který postavě dodá třeba nějaký bonus
nebo zvláštní schopnost. Definice může být precizní, ale týká se jenom hry a
souvislost se životem bude nulová. Navíc se tam pojem může zcela vyprázdnit.
V~rámci světa mohu svátost definovat vágně a napadnutelně např. jako viditelné
znamení Božího působení. Otevírají se zde dveře ke hnidopišství a otázkám, co
tedy ještě za svátost považovat a co nikoliv. Mohu vymýšlet různá viditelná
znamení, která by se za svátost považovala stěží, a definici tak nabourat.
Nicméně při dobré vůli rozumět se může takto definovaným pojmem předat něco
z~jeho vnitřního smyslu. Může pomoci k~pochopení a sladění myslí toho, kdo
definici tvoří či sděluje a toho, komu je předložena. Na třetím stupni, za
hranicemi světa, u Boha má svátost význam sama o sobě a její pochopení a prožití
je zázrak, který nejde mechanicky zprostředkovat. Odtud pojem pramení, jeho
prožití mu dává vzniknout v~lidském duchu a následně v~jazyce. Definovat ho tak,
aby došlo k~pochopení na této úrovni, je možné jen na základě skutečného
propojení dvou duší a užití slov či jiných prostředků přizpůsobených konkrétnímu
jedinci v~konkrétní situaci. Zde se pak můžeme bavit o pojmu zasvěcení spíše než
o běžném sdělování.

Jde-li nám tedy o to, ujasnit si, co míníme svátostí, nemůžeme toho na těchto
stránkách vyčerpávajícím způsobem dosáhnout. V~rámci modelu je to snadné: Pro
Církev československou husitskou je svátostí sedm, a to
\begin{enumerate}
\item{křest,}
\item{biřmování,}
\item{večeře Páně,}
\item{kněžské svěcení,}
\item{zpověď,}
\item{manželství a}
\item{útěcha nemocných.}
\end{enumerate}
Každá z~nich má pak samostatnou definici, a to opět na každé ze jmenovaných
úrovní.

V~rámci tohoto světa pojem svátosti definovali mnozí. K~nejčastěji zmiňovaným
patří sv. Augustýn: \textit{,,Sacramentum est sacrum signum,``} (De Civitate Dei
X, 5) čili ,,svátost je svaté znamení``. Známější verze
\textit{,,sacramentum est invisibilis gratiae visibilis forma,``} čili ,,svátost
je viditelná forma neviditelné milosti,`` je patrně Augustýnovi připisována
neprávem, byť harmonuje s~duchem jeho nauky \citep{king1967origin}.

Dalším velikánem je Tomáš Akvinský, který praví: \textit{,,Sacramentum potest
ali\-quid dici vel quia in se habet aliquam sanctitatem occultam, [...] vel quia
habet aliquem ordinem ad hanc sanctitatem,``} (Summa Theologiae III Q60) čili
,,Svátostí může být zváno něco, co v~sobě má nějakou skrytou svatost nebo co má
nějaký řád (rozuměj vztah) k~takové svatosti.``

Z~moderních autorů zmiňme Karla Bartha, který se svým způsobem drží Augustinova
pojetí: ,,Svátost je doprovodným a potvrzujícím symbolickým úkonem.``
\citep[I, 1]{barth1946kirchliche} Doprovodným se míní ke zvěstování, z~čehož pro
mne plyne pojetí svátosti jako čehosi principielně menšího než zvěstování.

K~doteku třetí, božské úrovně nám může dopomoci samotný dokument, kterým se
zaobírá tato práce. V~závěru úvodu stojí: \textit{,,Kéž tento naukový dokument
v~milosti a světle Ducha svatého je pomocí ke katechezi, vzdělávání ve víře,
zvěstování církve a jejímu liturgickému slavení a svátostné službě.``} Zde totiž
jazyk umožňuje dvojí čtení: svátostnou službu jednak jako vysluhování svátostí
definovaných na úrovni modelu, čili jedné ze sedmi, a jednak jako službu, která
je svým charakterem svátostná. Toto druhé, patrně nezamýšlené čtení může
evokovat libovolnou činnost, která je prováděna ve skutečné službě Duchu
svatému, a to ji činí svátostnou. Můžeme tak aplikovat Tomášův pojem svátosti
jako něčeho, co v~sobě skrytě obsahuje svatost. V~tomto duchu nechť je svátost
chápána a nechť je takto ve svátost proměněna veškerá naše činnost.

\vspace{15mm}
\chapter{Svátost v~Základech víry}
\label{div:zaklady}

\section{Mimo stať věnovanou svátostem}

Krom výše uvedeného užití obratu \textit{,,svátostná služba``} nacházíme zmínku
o svátosti v~citované pasáži \textbf{Čtyř pražských artikulů}. Ty jsou uvedeny jako jeden
z~pramenů našich Základů víry. Hovoří se zde o \textit{,,velebné svátosti těla a
krve Pána našeho Ježíše Krista,``}. Samo použití pozitivního elativu
\textit{velebný} má dalekosáhlé důsledky. Nastavuje to osobní a zúčastněný
postoj vůči svátosti, v~tomto případě večeře Páně. Je to svým způsobem vyznání a
bez něho se nelze svátosti přiblížit jinak než jako sterilnímu pojmu. Je proto
nanejvýš šťastné, že se Základy na tento starý text odvolávají jako na svůj
pramen.

Další užití slova svátost nacházíme v~odpovědi na \textbf{otázku 13}: \textit{,,v čem je
Boží církev po celém světě jednotná?{}``}, a sice že \textit{,,užívá svátosti křtu
a večeře Páně.``} Zdálo by se, že se vlastně zde o svátostech nic nevypovídá,
nýbrž že se vypovídá jen o církvi. Opak je pravdou: Svátosti křtu a večeře Páně
jsou zde postaveny do spojujícího prvku všech světových křesťanských denominací.
Dává jim to další rozměr posvátnosti: Krom inherentního významu tím, co jsou,
přibírají rozměr definujícího prvku celosvětového křesťanstva. ,,Jsem-li účasten
těchto dvou svátostí, patřím do rodiny křeťanů z~celé Země.``

V~odpovědi na \textbf{otázku 27} \textit{,,Kdo je hlavou Boží církve?{}``} se dočítáme, že
\textit{,,{},Hlavou` Boží církve je Kristus, který ji vede, oživuje a řídí svým
Duchem skrze svědectví Písma svatého a svátosti jako své tělo.``} Především je
zde třeba větu syntakticky zjednoznačnit. Nabízí se totiž čtení, že se hovoří o
\textit{svátostech jako (např.) Kristovo tělo}. To je čtení mylné a větu je
třeba číst tak, že \textit{Kristus církev řídí jako své tělo}, a to
\textit{skrze svědectví Písma svatého a svátosti}. Plyne to z~protikladu hlavy a
těla, kterého se zde užívá jako metafory.

Výpověď, že Kristus vede, oživuje a řídí církev skrze svátosti, je do jisté míry
problematická. Církev je vedena Kristem tehdy, když v~ní působí Duch svatý, to
považuji za natolik zřejmé, že pro to snad netřeba dále argumentovat. Výrok
v~odpovědi na otázku 27 by tedy implikoval, že působení Ducha svatého je
podmíněno vykonáváním svátostí. Toto připomíná mechanické středověké uvažování,
jako např. že kdo neprodělal církevní obřad křtu, nemůže vstoupit do nebe.

Je samozřejmě možné interpretovat výrok i tak, že vykonáváním a přijímáním
svátostí s~upřímností a pokorou se disponujeme přijetí Ducha svatého a ten nás
pak může vést, oživovat a řídit. Také lze výrok interpretovat tak, že se zde
svátost myslí obecnější, v~duchu Tomáše Akvinského jako to, co skrytě obsahuje
svatost. Potom k~tomu můžeme zajisté zařadit i přijetí Ducha svatého a říci, že
skrze takovou svátost Kristus církev řídí.

Vzhledem k~tomu, že však tento rozbor směřuje zevnitř církve v~jejím rámci k~ní
samotné, nepokládám za vhodné uplatňovat shovívavý postoj. Naopak, domnívám se,
že je vhodné poukázat na potenciální nedostatek -- ať už ve formulaci Základů
víry nebo, což je pravděpodobnější, v~mém vlastním jejich pochopení.

Navrhoval bych odpověď upravit například do této podoby: \textit{,,,Hlavou` Boží
církve je Kristus, který ji vede, oživuje a řídí svým Duchem, jehož nalézáme
v~Písmu svatém a o něhož se ucházíme přijímáním svátostí.``}

Odpověď na \textbf{otázku 93} zní: \textit{,,Zvěst církve je nutno doplňovat a potvrzovat
svátostmi a svědectvím života.``} Domnívám se, že v~tomto zdánlivě banálním
výroku se skrývá velká moudrost. Opět by sice bylo možné vidět zde mechanické
uvažování a upevňování církevní moci diktátem povinnosti přijímat svátosti. To
by však v~tomto případě ukazovalo spíše na takové uvažování na straně čtenáře,
než na chybu ve formulaci.

Nutnost svátostí je formulována i u Akvinského, opět v~Sumě III/60, ovšem to
neznamená, že by byla nějak samozřejmou. Konstatovat nutnost svátostí je zcela
na místě, a to proto, že jimi člověk rozšiřuje svůj život o posvátný rozměr.
Křtem zasvěcuje celý život Bohu. Bez křtu může žít pro sebe, jako pokřtěný už má
žít Bohu. Bez manželství může mít partnera pro svoje potěšení nebo výhody.
Manželstvím před Bohem se vztah stává prostředkem vztahu s~Bohem. Pokáním
opouští svoji hříšnost a zavazuje se zůstat čist. Přijímáním těla a krve Páně se
sjednocuje s~Bohem. Pochopitelně prázdné obřady nic z~toho nezprostředkují --
tolik se za sebe odvažuji tvrdit, ať je to v~rozporu s~jakoukoliv naukou. Ale
přijímá-li člověk svátosti opravdově, všechno toto se děje. A bez toho se to
děje stěží, protože naše životy nás táhnou k~žití pro sebe. Z~toho důvodu plně
souhlasím, že svátostmi a svědectvím života je nutno zvěst církve doplňovat.

\section{Část VI: O svátostech}

Konečně se dostáváme k~vlastní definici pojmu svátosti podle Základů víry. Podle
odpovědi na \textbf{otázku 307}
\textit{,,Svátost je jednání, jímž se obecenství věřících v~Duchu svatém činně
podílí na milosti Božího Slova přítomného skrze slyšené svědectví Písma svatého
a svátostné úkony.``} Tato formulace si zaslouží podrobnější rozbor.

\textit{,,Svátost je jednání.``} A sice jednání obecenství věřících. Na jednu
stranu je tato formulace kategoricky konkrétnější než u Augustina a Tomáše.
Nejde už o zcela neurčité \textit{znamení (,,signum``)} nebo \textit{něco
(,,aliquid``)}. Jde o jednání věřících. Předměty, jednání Boží, stavy mysli,
místa, slova, dokonce slova evangelia, vše je z~pojmu vyňato. Nabízí se otázka,
jak autoři formulace k~definici přistupovali: Zda jim šlo o definici
vyčerpávající nebo spíše deskriptivní. Jinými slovy zda jim šlo o to, vymezit,
co všechno lze nazvat svátostí a co už nikoliv, nebo zda jim šlo o to, úkony,
které zveme svátostmi co nejlépe charakterizovat. Vzhledem k~tomu, že se jedná o
odpověď na otázku \textit{,,Co je svátost?``}, logicky bych čekal definici
vyčerpávající, tedy formulaci, která vymezí, co svátost je a co svátost není.
Ze způsobu odpovědi však usuzuji, že opdověď spíše charakterizuje ty úkony,
které zveme svátostmi. Každopádně potom je odpověď mnohem legitimnější než
v~případě vyčerpávající definice.

Pro demonstraci toho, že jako vyčerpávající definice formulace ze Základů víry
nepostačuje, si vezměme příklad mezního křtu: Dva nevěřící lidé Alice a
Bořek\footnote{Podle vzoru scénářů bezpečnostních rizik pojmenovávám postavy
jmény podle abecedy. V~angličtině se standardně používá Alice a Bob, viz
\url{https://en.wikipedia.org/wiki/Alice\_and\_Bob}.} jsou spolu sami a Alice je
v~ohrožení života. Alice si vzpomene na to, že babička kladla důraz na
náboženství a zatouží po křtu. Bořek křtu nepřikládá žádný význam, ale ze sympatie
Alici polije vodou a pronese křestní formuli. Alici se přitíží a za několik
minut se dostane na pokraj smrti. Tehdy se jí dostane mystická zkušenost a stane
se z~ní hluboce věřící člověk.

Jedná se zde o svátost nebo nikoliv? Podle zdravého úsudku jistě ano: Křest byl
církevně přípustný a vyslyšený Bohem. Podle definice ale nikoliv, protože nešlo
o jednání obecenství věřících: Během křtu byli oba dva nevěřící.
Přikloňme se tedy k~variantě, že formulace má popisovat to, co v~církvi zveme
svátostí. Vhodnější otázkou by tedy mohlo být: \textit{,,Čím jsou svátosti?{}``}

\textit{,,... podílí na milosti Božího Slova...``} Co znamená podílet se na
milosti Božího slova? Je zde míněno Boží slovo ve smyslu biblického textu nebo
ve smyslu otázky 52 jako Bůh sám? Pokud se míní Bůh sám, proč se užívá pojmu
Boží slovo? Co se míní milostí Božího slova? Je to milost, kterou udílí Boží
slovo? Nebo je to milost, ze které nám Boží slovo je dáno? V~jakém smyslu se
míní podílení? Jde o to, že my spoluudílíme milost s~Bohem? Nebo že jsme jedni
z~těch, kterým je milost udělena (máme na ní podíl)?

\textit{,,...v~Duchu svatém činně podílí...``} Podílet se v~Duchu svatém a
podílet se činně jsou dvě příslovečná určení. Mají deskriptivní nebo
restriktivní význam? To jest: Vypovídá se zde, že podílení je v~Duchu svatém a
je činné, nebo to znamená, že jiného podílení než činného a v~Duchu svatém se
tato věta netýká?

Abych rozdíl ilustroval, vezměme si příklad věty: 

Už jen v~této samotné slovní vazbě můžeme prostým kartézským součinem dospět
k~osmi možným variantám. Všechny je rozebrat je holý nesmysl. Pokusme se proto o
intuitivní zjednoznačnění -- jak byla formulace asi myšlena: Při čtení a kázání
se předává svědectví Písma svatého. Také se provádějí svátostné úkony. V~těchto
dvou formách je přítomné Boží slovo. Božím slovem se míní Bůh, který se
v~biblickém textu zjevuje. Slovo Boží je jednak předmětem milosti, tedy
z~milosti Bohem udělené je nám Boží slovo dáno, jednak původcem milosti, tedy
Boží slovo coby Bůh udílí milost. Víceznačnost ponecháváme a rozumíme oběma
způsoby. Obecenství věřících slyší čtení, kázání a účastní se svátostných úkonů.

\vspace{15mm}
\chapter{Závěr}
\label{div:zaver}



%%% Seznam použité literatury
%%% Seznam použité literatury (bibliografie)
%%%
%%% Pro vytváření bibliografie používáme bibTeX. Ten zpracovává
%%% citace v textu (např. makro \cite{...}) a vyhledává k nim literaturu
%%% v souboru literatura.bib.
%%%
%%% Příkaz \bibliographystyle určuje, jakým stylem budou citovány odkazy
%%% v textu. V závorce je název zvoleného souboru .bst. Styly plainnat
%%% a unsrt jsou standardní součástí latexových distribucí. Styl czplainnat
%%% je dodáván s touto šablonou a bibTeX ho hledá v aktuálním adresáři.

\bibliographystyle{czplainnat}    %% Autor (rok) s českými spojkami
% \bibliographystyle{plainnat}    %% Autor (rok) s anglickými spojkami
% \bibliographystyle{unsrt}       %% [číslo]

\renewcommand{\bibname}{Seznam použité literatury}

%%% Vytvoření seznamu literatury. Pozor, pokud jste necitovali ani jednu
%%% položku, seznam se automaticky vynechá.

\bibliography{citace}

%%% Kdybyste chtěli bibliografii vytvářet ručně (bez bibTeXu), lze to udělat
%%% následovně. V takovém případě se řiďte normou ISO 690 a zvyklostmi v oboru.

% \begin{thebibliography}{99}
%
% \bibitem{lamport94}
%   {\sc Lamport,} Leslie.
%   \emph{\LaTeX: A Document Preparation System}.
%   2. vydání.
%   Massachusetts: Addison Wesley, 1994.
%   ISBN 0-201-52983-1.
%
% \end{thebibliography}


%%% Obrázky v disertační práci
%%% (pokud jich je malé množství, obvykle není třeba seznam uvádět)
%\listoffigures

%%% Tabulky v disertační práci (opět nemusí být nutné uvádět)
%%% U matematických prací může být lepší přemístit seznam tabulek na začátek práce.
%\listoftables

%%% Použité zkratky v disertační práci (opět nemusí být nutné uvádět)
%%% U matematických prací může být lepší přemístit seznam zkratek na začátek práce.
%\chapwithtoc{Seznam použitých zkratek}

%%% Součástí doktorských prací musí být seznam vlastních publikací
%\input{publikace.tex}

%%% Přílohy k disertační práci, existují-li. Každá příloha musí být alespoň jednou
%%% odkazována z vlastního textu práce. Přílohy se číslují.
%%%
%%% Do tištěné verze se spíše hodí přílohy, které lze číst a prohlížet (dodatečné
%%% tabulky a grafy, různé textové doplňky, ukázky výstupů z počítačových programů,
%%% apod.). Do elektronické verze se hodí přílohy, které budou spíše používány
%%% v elektronické podobě než čteny (zdrojové kódy programů, datové soubory,
%%% interaktivní grafy apod.). Elektronické přílohy se nahrávají do SISu a lze
%%% je také do práce vložit na CD/DVD. Povolené formáty souborů specifikuje
%%% opatření rektora č. 72/2017.
%\appendix
%\chapter{Přílohy}

%\section{První příloha}

\openright
\end{document}
